% GENERATED FILE. Edit canonical lessons, then run npm run generate:books.
% canonical-source-hash: fnv1a64:f6faa2b884fadb4d
% canonical-lessons: FA-C466-galle, FA-C466-chupan, FA-C466-tavile, FA-C466-marta, FA-C466-charidan

\chapter{Animals in the fields}
\label{ch:fa-animals-in-the-fields}
\hlchaptermodality{466}

\begin{chapteropening}
\textbf{By the end of this chapter:} \emph{I can talk about animals in the countryside: a flock, a shepherd, a stable, a pasture and grazing.}

Complete the last lesson of chapter 466: I can talk about animals in the countryside: a flock, a shepherd, a stable, a pasture and grazing.
\end{chapteropening}

\section[\texorpdfstring{\fa{گله} (galle)}{گله (galle)}]{\fa{گله} (galle) — a flock, a herd}

\label{lesson:FA-C466-galle}



\begin{quote}
Before the new one: say the Persian for a vineyard, then the Persian for to plough.
\end{quote}

\subsection*{You'll want to know: \fa{گله}}
\textbf{\fa{گله}} (\emph{galle}) — \textquotedblleft{}a flock, a herd\textquotedblright{}.

\begin{quote}
\fa{یک} \fa{گله} \fa{گوسفند} \fa{از} \fa{جاده} \fa{رد} \fa{شد}.
\end{quote}

\begin{quote}
\emph{Yek galle gusfand az jâdde rad shod.}
\end{quote}

\begin{quote}
A flock of sheep crossed the road.
\end{quote}

Do not mix it up with \emph{gele} 'a complaint', written the same way.

\begin{cousinweb}
\emph{galle} 'herd' is a native word whose older form is not settled. EtymDB lists only its look-alike \emph{gele} 'complaint'.
\end{cousinweb}

\begin{grammarlens}[title={What you've built}]
One more word for telling, explaining and describing at B1.
\end{grammarlens}

\subsection*{Your turn}
\begin{itemize}
\raggedright
  \item \emph{Say it:} \emph{galle}
  \item \emph{Say it:} \emph{galle}, once more
  \item \emph{Say it:} a sentence of your own with \emph{galle}
  \item \emph{Recall:} say the Persian for a vineyard, then the Persian for to plough, then say \emph{galle} again
\end{itemize}

\begin{tcolorbox}[breakable,colback=teal!4,colframe=teal!35!black,title={Before you move on}]
What does \fa{گله} mean? (\textquotedblleft{}A flock, a herd\textquotedblright{}.) Say it once more.
\end{tcolorbox}
\section[\texorpdfstring{\fa{چوپان} (chupân)}{چوپان (chupân)}]{\fa{چوپان} (chupân) — a shepherd}

\label{lesson:FA-C466-chupan}



\begin{quote}
Before the new one: say the Persian for to plough, then the Persian for a flock.
\end{quote}

\subsection*{You'll want to know: \fa{چوپان}}
\textbf{\fa{چوپان}} (\emph{chupân}) — \textquotedblleft{}a shepherd\textquotedblright{}.

\begin{quote}
\fa{چوپان} \fa{گوسفندها} \fa{را} \fa{به} \fa{کوه} \fa{برد}.
\end{quote}

\begin{quote}
\emph{Chupân gusfand-hâ râ be kuh bord.}
\end{quote}

\begin{quote}
The shepherd took the sheep to the mountain.
\end{quote}

Turkish \emph{çoban} is the same word.

\begin{cousinweb}
\emph{chupân} is a native word, which EtymDB lists without older forms; it is usually traced to Middle Persian \emph{šubān} 'shepherd'.
\end{cousinweb}

\begin{grammarlens}[title={What you've built}]
One more word for telling, explaining and describing at B1.
\end{grammarlens}

\subsection*{Your turn}
\begin{itemize}
\raggedright
  \item \emph{Say it:} \emph{chupân}
  \item \emph{Say it:} \emph{chupân}, once more
  \item \emph{Say it:} a sentence of your own with \emph{chupân}
  \item \emph{Recall:} say the Persian for to plough, then the Persian for a flock, then say \emph{chupân} again
\end{itemize}

\begin{tcolorbox}[breakable,colback=teal!4,colframe=teal!35!black,title={Before you move on}]
What does \fa{چوپان} mean? (\textquotedblleft{}A shepherd\textquotedblright{}.) Say it once more.
\end{tcolorbox}
\section[\texorpdfstring{\fa{طویله} (tavile)}{طویله (tavile)}]{\fa{طویله} (tavile) — a stable, a cowshed}

\label{lesson:FA-C466-tavile}



\begin{quote}
Before the new one: say the Persian for a flock, then the Persian for a shepherd.
\end{quote}

\subsection*{You'll want to know: \fa{طویله}}
\textbf{\fa{طویله}} (\emph{tavile}) — \textquotedblleft{}a stable, a cowshed\textquotedblright{}.

\begin{quote}
\fa{اسبها} \fa{شب} \fa{در} \fa{طویله} \fa{میخوابند}.
\end{quote}

\begin{quote}
\emph{Asb-hâ shab dar tavile mikhâband.}
\end{quote}

\begin{quote}
The horses sleep in the stable at night.
\end{quote}

For sheep the word is \emph{âghol}.

\begin{cousinweb}
\emph{tavile} was borrowed from Arabic and is usually traced to Arabic \emph{ṭawīla} 'a long rope for tethering animals'. EtymDB does not list it.
\end{cousinweb}

\begin{grammarlens}[title={What you've built}]
One more word for telling, explaining and describing at B1.
\end{grammarlens}

\subsection*{Your turn}
\begin{itemize}
\raggedright
  \item \emph{Say it:} \emph{tavile}
  \item \emph{Say it:} \emph{tavile}, once more
  \item \emph{Say it:} a sentence of your own with \emph{tavile}
  \item \emph{Recall:} say the Persian for a flock, then the Persian for a shepherd, then say \emph{tavile} again
\end{itemize}

\begin{tcolorbox}[breakable,colback=teal!4,colframe=teal!35!black,title={Before you move on}]
What does \fa{طویله} mean? (\textquotedblleft{}A stable, a cowshed\textquotedblright{}.) Say it once more.
\end{tcolorbox}
\section[\texorpdfstring{\fa{مرتع} (marta')}{مرتع (marta')}]{\fa{مرتع} (marta') — a pasture, grazing land}

\label{lesson:FA-C466-marta}



\begin{quote}
Before the new one: say the Persian for a shepherd, then the Persian for a stable.
\end{quote}

\subsection*{You'll want to know: \fa{مرتع}}
\textbf{\fa{مرتع}} (\emph{marta'}) — \textquotedblleft{}a pasture, grazing land\textquotedblright{}.

\begin{quote}
\fa{در} \fa{بهار} \fa{مرتع} \fa{سبز} \fa{است}.
\end{quote}

\begin{quote}
\emph{Dar bahâr marta' sabz ast.}
\end{quote}

\begin{quote}
In spring the pasture is green.
\end{quote}

Arabic plural \emph{marâte'}.

\begin{cousinweb}
\emph{marta'} was borrowed from Arabic \emph{martaʿ} 'pasture', which EtymDB derives from \emph{ratʿ} 'grazing'.
\end{cousinweb}

\begin{grammarlens}[title={What you've built}]
One more word for telling, explaining and describing at B1.
\end{grammarlens}

\subsection*{Your turn}
\begin{itemize}
\raggedright
  \item \emph{Say it:} \emph{marta'}
  \item \emph{Say it:} \emph{marta'}, once more
  \item \emph{Say it:} a sentence of your own with \emph{marta'}
  \item \emph{Recall:} say the Persian for a shepherd, then the Persian for a stable, then say \emph{marta'} again
\end{itemize}

\begin{tcolorbox}[breakable,colback=teal!4,colframe=teal!35!black,title={Before you move on}]
What does \fa{مرتع} mean? (\textquotedblleft{}A pasture, grazing land\textquotedblright{}.) Say it once more.
\end{tcolorbox}
\section[\texorpdfstring{\fa{چریدن} (charidan)}{چریدن (charidan)}]{\fa{چریدن} (charidan) — to graze}

\label{lesson:FA-C466-charidan}



\begin{quote}
Before the new one: say the Persian for a stable, then the Persian for a pasture.
\end{quote}

\subsection*{You'll want to know: \fa{چریدن}}
\textbf{\fa{چریدن}} (\emph{charidan}) — \textquotedblleft{}to graze\textquotedblright{}.

\begin{quote}
\fa{گاوها} \fa{در} \fa{دشت} \fa{میچرند}.
\end{quote}

\begin{quote}
\emph{Gâv-hâ dar dasht micharand.}
\end{quote}

\begin{quote}
The cows are grazing on the plain.
\end{quote}

Present stem \emph{char-}: \emph{micharad}. \emph{Charândan} is to graze a flock.

\begin{cousinweb}
\emph{charidan} is inherited: EtymDB traces it to the PIE root \emph{k\textsuperscript{w}el-} 'to move, to turn (around)', as animals wander while they graze.
\end{cousinweb}

\begin{grammarlens}[title={What you've built}]
One more word for telling, explaining and describing at B1.
\end{grammarlens}

\subsection*{Your turn}
\begin{itemize}
\raggedright
  \item \emph{Say it:} \emph{charidan}
  \item \emph{Say it:} \emph{charidan}, once more
  \item \emph{Say it:} a sentence of your own with \emph{charidan}
  \item \emph{Recall:} say the Persian for a stable, then the Persian for a pasture, then say \emph{charidan} again
\end{itemize}

\begin{tcolorbox}[breakable,colback=teal!4,colframe=teal!35!black,title={Before you move on}]
What does \fa{چریدن} mean? (\textquotedblleft{}To graze\textquotedblright{}.) Say it once more.
\end{tcolorbox}
